% SOURCE_AUTHORITY: sga5_fr_workpass.tex, lines 10587--10609 (LF-normalized unit).
% SOURCE_SHA256: 791F4EFFC5E02832D5D77ED03518C8156D6F07E4C8238B03545DB93D883FBB28
\subsection*{7.7}

Por analogia com o que ocorre no caso transcendente, onde para coeficientes constantes se dispõe de outro teorema bem conhecido de Lefschetz ([4]), chega-se a conjecturar o enunciado seguinte. Se $X$ é um $k$-esquema projetivo, liso e conexo de dimensão $n$, e $\xi$ é a classe em $H^2(X,\mathbb Z_\ell(1))$ de um $\mathcal O_X$-módulo inversível muito amplo, então, para todo $\mathbb Q_\ell$-feixe constante torcido $L$ e todo inteiro $p\le n$, a multiplicação por $\xi^{n-p}$:
\[
H^p(X,L)\to H^{2n-p}(X,L(n-p))
\]
é um isomorfismo. Como observou Deligne, os enunciados (7.7) são falsos sem uma hipótese suplementar sobre $L$. No entanto, Deligne os demonstrou (A conjectura de Weil II) quando, além disso, $L$ é puro.

Este enunciado conjectural tem as consequências seguintes, tampouco demonstradas em $\operatorname{car}\ne0$.

Sejam $X$ um $k$-esquema projetivo liso, conexo, de dimensão $n$, $Y$ uma seção hiperplana lisa de $X$, e $L$ um $\mathbb Q_\ell$-feixe constante torcido sobre $X$.
\begin{enumerate}[label=(\roman*)]
\item O morfismo de imagem inversa
\[
H^p(X,L)\to H^p(Y,L|Y)
\]
é um epimorfismo para $p\ge n$.
\item O morfismo de Gysin
\[
H^p(Y,(L|Y)(-1))\to H^{p+2}(X,L)
\]
é um monomorfismo para $p\le n-2$.
\end{enumerate}
